%! Exponential Function Context and Data Modeling — Unit 4, Lesson 4.6 — Lesson Plan
\documentclass[10pt]{article}
\usepackage{precalculus-article}
\usepackage{precalculus-boxes}
\usepackage{graphicx}
\graphicspath{{images/}}
\newcommand{\TallMath}[1]{$\displaystyle #1\rule[-1.4em]{0pt}{3.2em}$}
\newcommand{\CourseName}{Precalculus}
\newcommand{\SchoolYear}{2026--2027}
\newcommand{\MeetingLength}{55 minutes}
\newcommand{\UnitNumberName}{Unit 4: Exponential Functions \quad}
\newcommand{\LessonNumberName}{Lesson 4.6: Exponential Function Context and Data Modeling}

\begin{document}
\begin{center}
    {\Large\bfseries \CourseName: \SchoolYear \\
    {\small\bfseries \UnitNumberName \LessonNumberName}}
\end{center}

\begin{tcolorbox}[colback=sky, colframe=navy, boxrule=0.9pt, arc=2mm,
    left=2mm, right=2mm, top=1.5mm, bottom=1.5mm]
    \textbf{Primary Objective:} Students will construct exponential function models from
    contextual data (using a growth/decay ratio and initial value, or two input-output pairs),
    interpret the base as a growth factor and percent change, and apply models to make
    predictions within a reasonable domain.
    \hfill\textit{\small Skills: 1.C and 3.B}
\end{tcolorbox}

\renewcommand{\arraystretch}{1.6}
\begin{skillbox}[Priority Ideas \& Skills]{goldbox}
\begin{tabularx}{\linewidth}{@{}X X@{}}
\textbf{AP Skills} &
\textbf{Key Understandings} \\
\midrule
\textbf{Skill 1.C} --- Identify mathematical information from graphical, numerical, analytical, and verbal representations. &
\begin{itemize}[leftmargin=1.2em,itemsep=2pt,topsep=0pt]
  \item Successive output ratios over equal-length input intervals reveal exponential structure. (EK 2.5.A.1)
  \item A model $f(x) = ab^x$ is built from an initial value $a$ and base $b$, or solved from two points. (EK 2.5.A.3)
  \item The natural base $e \approx 2.718$ appears in continuous growth models. (EK 2.5.A.6)
\end{itemize} \\
\textbf{Skill 3.B} --- Apply an appropriate mathematical definition, theorem, or test. &
\begin{itemize}[leftmargin=1.2em,itemsep=2pt,topsep=0pt]
  \item Base $b$ is the growth factor; $b = 1+r$ for growth, $b = 1-r$ for decay. (EK 2.5.B.1)
  \item Equivalent forms reveal different time-scale properties: $2^d = (2^7)^{d/7}$. (EK 2.5.B.2)
  \item Models predict values within contextual domain constraints. (EK 2.5.B.3)
\end{itemize} \\
\end{tabularx}
\end{skillbox}

\begin{skillbox}[{Vocabulary, Concepts, \& Theorems}]{greenbox}
\renewcommand{\arraystretch}{1.4}
\begin{tabularx}{\linewidth}{@{} >{\bfseries\color{navy}}p{0.27\linewidth} X @{}}
exponential model      & A function $f(x) = ab^x$ that models proportional growth or decay over equal-length intervals. \\
growth factor          & The base $b$ in $f(x) = ab^x$; the multiplier applied at each unit step. \\
percent change         & The rate $r$ such that $b = 1+r$ (growth) or $b = 1-r$ (decay); expressed as a percentage. \\
initial value          & The value $a = f(0)$; the output when the input equals zero. \\
natural base $e$       & The irrational number $e \approx 2.718$; used as the base in continuous growth/decay models. \\
exponential regression & A technology-produced best-fit exponential model for a data set. \\
prediction             & Using the model to estimate an output for an input not in the original data. \\
domain restriction     & Limiting the input values of a model to those that are meaningful in context. \\
\end{tabularx}
\end{skillbox}

\begin{skillbox}[Activate Prior Knowledge \& Spiral Review]{sky}
    Spiral review (warm-up) covers skills from Lesson 2.4: rewriting exponential expressions
    using exponent rules. Students rewrite $f(x) = 2^{x+5}$ in $a \cdot 2^x$ form, convert
    $4^{3x}$ to a single base to the $x$ power, and simplify $9^{3/2}$.
\end{skillbox}

\begin{skillbox}[Hook]{sky}
Present the following population data on the board:

\medskip
\begin{center}
\begin{tabular}{c|c}
\textbf{Year} & \textbf{Population} \\ \hline
0 & 10{,}000 \\
1 & 10{,}300 \\
2 & 10{,}609 \\
3 & 10{,}927 \\
\end{tabular}
\end{center}

\medskip
Ask: \textbf{``Is this growth linear or exponential? What is the growth rate per year?''}
After student responses, compute successive ratios: $10300/10000 = 1.03$, $10609/10300 \approx 1.03$,
etc. Reveal: \textit{equal ratios over equal intervals --- this is exponential with a 3\% annual
growth rate.} Today's lesson: how do we build and use such a model?
\end{skillbox}

\begin{skillbox}[Lesson]{sky}
\begin{multicols}{2}
\textbf{Part 1 (15 min): Build from Ratio and Initial Value.}

If $b$ is the growth factor and $a$ is the initial value:
$f(x) = a \cdot b^x$. For the hook: $f(t) = 10000 \cdot (1.03)^t$.

Connect percent change to base: growth at rate $r \Rightarrow b = 1+r$;
decay at rate $r \Rightarrow b = 1-r$.

Model 3 examples: (1) 5\% annual growth from 200; (2) 12\% annual decay from 500;
(3) given a table --- check ratios, read off $a$ and $b$.

\columnbreak

\textbf{Part 2 (15 min): Build from Two Points.}

Given $(x_1, y_1)$ and $(x_2, y_2)$, set up the system:
$y_1 = ab^{x_1}$ and $y_2 = ab^{x_2}$.
Divide to eliminate $a$: $b^{x_2 - x_1} = y_2/y_1$; solve for $b$, then $a$.

Model full example: population 4000 at year 2, 6250 at year 5 (Activity Tier A preview).

Introduce the natural base $e \approx 2.718$: used in continuous models, $f(t) = ae^{rt}$.
Compare to discrete base: $e^{0.03} \approx 1.0305$ (≈ 3.05\% continuous growth).

\end{multicols}

\medskip
\textbf{Part 3 (10 min): Apply the Model --- Prediction and Equivalent Forms.}

Emphasize domain cautions: does the model apply far in the future?

Equivalent forms (EK 2.5.B.2): $f(d) = 3 \cdot 2^d$ models daily growth. Rewrite for weekly:
$f(d) = 3 \cdot (2^7)^{d/7} = 3 \cdot 128^{d/7}$. Each unit of $d/7$ is one week, factor is 128.
What does $128$ tell us? The quantity multiplies by 128 each week.
\end{skillbox}

\begin{skillbox}[Active Monitoring]{redbox}
\begin{itemize}[leftmargin=1.2em,itemsep=2pt,topsep=0pt]
  \item After Part 1: cold-call ``If the base is 0.88, is this growth or decay? What is the percent change?''
        Expect: decay, 12\% per period.
  \item During two-point model: watch for students who divide equations incorrectly or forget to
        take the correct root when solving for $b$.
  \item After Part 2: cold-call ``Why can we divide the two equations to eliminate $a$?''
        Expect: $a$ is a common factor, dividing creates a ratio equation with only $b$.
  \item During Part 3: confirm students understand that the base changes when the unit of time changes
        but the underlying growth is the same function.
\end{itemize}
\end{skillbox}

\begin{skillbox}[Group Work \& Differentiation]{redbox}
\begin{itemize}[leftmargin=1.2em,itemsep=2pt,topsep=0pt]
  \item \textbf{Tier R (Remediate):} Drug concentration model with half-life structure given; build
        $C(t) = 200 \cdot (0.5)^{t/4}$; complete a table; interpret the base; predict one value.
  \item \textbf{Tier A (Apply):} Two data points given; construct the model algebraically; find the
        annual growth rate as a percent; predict a future value with domain note.
  \item \textbf{Tier E (Extend):} Rewrite a daily model as weekly and 3-day models; explain what each
        base means; determine the input when the quantity exceeds 10{,}000.
\end{itemize}
\end{skillbox}

\begin{skillbox}[Individual Work \& Assessment]{redbox}
Exit ticket (3 questions, 1 page): (1) write $P(t) = 5000 \cdot (1.06)^t$ and evaluate $P(10)$;
(2) find $a$ and $b$ from two points $(1, 6)$ and $(3, 54)$;
(3) interpret what $b = 1.06$ means in context.

Sort by: students who miss item 1 need review of building from percent change (LO 2.5.A);
students who miss item 2 need re-teaching of the two-point system (LO 2.5.A.3);
students who miss item 3 need work on contextual interpretation of $b$ (LO 2.5.B.1).
\end{skillbox}

\begin{skillbox}[Reinforcement \& Extension]{goldbox}
Homework (6 problems): convert 8\% annual decay to $ab^t$ form; two-point model with full
algebraic work; radioactive decay half-life; interpret equivalent forms in different time units;
AP-style MC with $f(t) = 500e^{0.04t}$; extension with a vertical shift revealing exponential pattern.

Extension: when data has a horizontal asymptote (e.g., a cooling problem), adding a constant to
the dependent variable values can reveal the proportional growth pattern --- preview of transformed
exponential models.

Preview of Lesson 2.6: logarithmic functions as inverses of exponentials --- how to ``undo'' an
exponential model to solve for the input (e.g., when will the population reach 20{,}000?).
\end{skillbox}

\end{document}
